% Part: computability
% Chapter: computability-theory
% Section: prop-reduce

\documentclass[../../../include/open-logic-section]{subfiles}

\begin{document}

\olfileid[hi]{cmp}{thy}{ppr}
\olsection{अपचयनीयता के गुण}

हम $A \leq_m B$ लिखते हैं यदि $A$,~$B$ में अपचयित होता है। यह संकेत
सुझाता है कि यदि $A \leq_m B$, तो $A$,~$B$ से ``अधिक कठिन नहीं'' और $B$,
~$A$ के ``बराबर या उससे अधिक कठिन'' है; तथा $\le_m$, संख्याओं पर सामान्य
$\le$ की तरह, समुच्चयों (या निर्णय समस्याओं) को उनकी जटिलता के अनुसार
क्रमबद्ध करता है। अगले दो साध्य इस सहज बोध का समर्थन करते हैं। पहला कहता
है कि $\le_m$ संक्रामक है।

\begin{prop}
\ollabel{prop:trans-red}
यदि $A \leq_m B$ और $B \leq_m C$, तो $A \leq_m C$।
\end{prop}

\begin{proof}
$A$ से~$B$ में अपचयन को $B$ से $C$ में अपचयन के साथ संयोजित करने पर
$A$ से~$C$ में अपचयन मिलता है।
\end{proof}

\begin{prob}
यह दिखाकर \olref{prop:trans-red} सिद्ध करें कि यदि $f$ और~$g$ क्रमशः
$A$ से~$B$ तथा $B$ से~$C$ में अनेक-एक अपचयन हैं, तो $\comp{f}{g}$,
$A$ से~$C$ में अनेक-एक अपचयन है।
\end{prob}

\begin{prop}
\ollabel{prop:reduce}
$A$ और $B$ को कोई समुच्चय मानें और मान लें कि $A \leq_m B$।
\begin{enumerate}
\item यदि $B$ संगणनीयतः परिगणनीय है, तो~$A$ भी है।
\item यदि $B$ संगणनीय है, तो~$A$ भी है।
\end{enumerate}
\end{prop}

\begin{proof}
$f$ को $A$ से~$B$ में अनेक-एक अपचयन मानें। पहले दावे के लिए केवल जाँचें
कि यदि $B$ किसी आंशिक फलन~$g$ का प्रांत है, तो $A$,~$\comp{f}{g}$ का
प्रांत है:
\begin{align*}
x \in A & \text{ तभी और केवल तभी } f(x) \in B \\
& \text{ तभी और केवल तभी }  g(f(x)) \fdefined.
\end{align*}

दूसरे दावे के लिए याद करें कि यदि $B$~संगणनीय है, तो $B$ और
~$\Complement{B}$ संगणनीयतः परिगणनीय हैं
(\olref[cmp]{thm:ce-comp})। यह जाँचना कठिन नहीं कि $f$,~$\Complement{A}$
से $\Complement{B}$ में भी अनेक-एक अपचयन है; अतः इस प्रमाण के पहले भाग
से $A$ और~$\Complement{A}$ दोनों !!{computably
enumerable} हैं। इसलिए
\olref[cmp]{thm:ce-comp} से $A$ भी संगणनीय है। (वैकल्पिक रूप से, आप जाँच
सकते हैं कि $\Char{A} = \comp{f}{\Char{B}}$; अतः यदि $\Char{B}$ संगणनीय
है, तो~$\Char{A}$ भी है।)
\end{proof}

\begin{prob}
मान लें कि $f$, $A$ से~$B$ में अनेक-एक अपचयन है। दिखाएँ कि $f$,
$\Complement{A}$ से $\Complement{B}$ में भी अनेक-एक अपचयन है।
\end{prob}

\begin{prob}
दिखाएँ कि यदि $f\colon A \to B$ अनेक-एक अपचयन है, तो $\Char{A}
= \comp{f}{\Char{B}}$।
\end{prob}

\begin{digress}
अपचयनीयता की अधिक सामान्य धारणा, \emph{ट्यूरिंग अपचयनीयता}, अन्य
संदर्भों में—विशेषतः अनिर्णेयता के परिणाम सिद्ध करने में—उपयोगी है।
ध्यान दें कि \olref[cmp]{cor:comp-k} से~$K_0$ का पूरक~$K_0$ में अपचयनीय
नहीं है, क्योंकि वह संगणनीयतः परिगणनीय नहीं है। किंतु सहज रूप से, यदि
आप $K_0$ के बारे में प्रश्नों के उत्तर जानते, तो उसके पूरक के बारे में
प्रश्नों के उत्तर भी जानते। कोई समुच्चय $A$,~$B$ में ट्यूरिंग अपचयनीय
कहलाता है यदि~$A$ के प्रश्नों के उत्तर,~$B$ के बारे में प्रश्न पूछ सकने
वाली संगणनीय प्रक्रिया से निर्धारित किए जा सकें। यह अनेक-एक अपचयनीयता
से अधिक उदार है, जिसमें (1)~आपको $B$ के बारे में केवल एक प्रश्न पूछने
की अनुमति है, और (2)~``हाँ'' उत्तर का $A$ के प्रश्न के ``हाँ'' उत्तर में
बदलना आवश्यक है, तथा ``नहीं'' के लिए भी वैसा ही। फिर भी, यदि $A$,~$B$
में ट्यूरिंग अपचयनीय और $B$ संगणनीय है, तो $A$ भी संगणनीय है (यद्यपि,
जैसा हमने देखा, संगणनीय परिगणनीयता के लिए अनुरूप कथन सत्य नहीं है)।

हमने जिन विभिन्न अपचयनीयता-धारणाओं पर चर्चा की है, उन पर विचार करें और
उनके भेद समझें। किंतु इस अध्याय में हम केवल अनेक-एक अपचयनीयता से व्यवहार
करेंगे। प्रसंगवश, पिछले अनुच्छेद में चर्चित दोनों प्रकारों के संगणनात्मक
जटिलता में अनुरूप रूप हैं, साथ में यह अतिरिक्त आवश्यकता कि ट्यूरिंग
मशीनें बहुपद समय में चलें: अनेक-एक अपचयनीयता का जटिलता-संस्करण
\emph{कार्प अपचयनीयता} और ट्यूरिंग अपचयनीयता का जटिलता-संस्करण
\emph{कुक अपचयनीयता} कहलाता है।
\end{digress}

\end{document}
